% chktex-file 1
% chktex-file 13

\chapter{Running Docker Locally}
\label{S:docker-local}
\index{Docker!Local}

\FILENAME\

The official installation documentation for docker can be found by
visiting the following Web page:

\URL{https://www.docker.com/community-edition}

Here you will find a variety of packages, one of which will hopefully
suitable for your operating system. The supported operating systems
currently include:

\begin{itemize}
\item  OSX, Windows, Centos, Debian, Fedora, Ubuntu, AWS, Azure
\end{itemize}

Please chose the one most suitable for you. For your convenience we
provide you with installation instructions for OSX
(Section~\ref{s:docker-osx}), Windows 10 (Section~\ref{s:docker-windows})
and Ubuntu (Section~\ref{s:docker-ubuntu}).

\section{Instillation for OSX}
\label{s:docker-osx}
\index{Docker!Install!OSX}

The docker community edition for OSX can be found at the following link

\URL{https://store.docker.com/editions/community/docker-ce-desktop-mac}

We recommend that at this time you get the version \textit{Docker CE for MAC (stable)}

\URL{https://download.docker.com/mac/stable/Docker.dmg}

Clicking on the link will download a dmg file to your machine, that
you than will need to install by double clicking and allowing access
to the dmg file. Upon installation a \texttt{whale} in the top status
bar shows that Docker is running, and you can access it via a terminal.

\begin{figure}[htb]
\centering
\includegraphics[width=0.5\textwidth]{whale-in-menu-bar.png}
\caption{Docker integrated in the menu bar on OSX}
\end{figure}

\section{Installation for Ubuntu}
\label{s:docker-ubuntu}
\index{Docker!Install!Ubuntu}

In order to instal Docker community edition for Ubuntu, you first have
to register the repository from where you can download it. This can be
achieved as follows:

\begin{lstlisting}[language=bash]
$ sudo apt-get update
$ sudo apt-get install \
    apt-transport-https \
    ca-certificates \
    curl \
    software-properties-common
$ curl -fsSL https://download.docker.com/linux/ubuntu/gpg | sudo apt-key add -
$ sudo apt-key fingerprint 0EBFCD88
$ sudo add-apt-repository \
   "deb [arch=amd64] https://download.docker.com/linux/ubuntu \
   $(lsb_release -cs) \
   stable"
\end{lstlisting}

Now that you have configured the repository location, you can install it
after you have updated the operating system. The update and install is
done as follows:

\begin{lstlisting}[language=bash]
$ sudo apt-get update
$ sudo apt-get install docker-ce
$ sudo apt-get update
\end{lstlisting}
%$


\section{Draft: Installation for Windows 10}
\label{s:docker-windows}
\index{Docker!Install!Windows}

\begin{WARNING}
The instructions for Windows have not passed our quality control. We
look for TAs and student that test out and improve this section.
\end{WARNING}

\TODO{TA: Docker on windows. Please improve and finalize the section}

Before we start we create the following directory:

\begin{lstlisting}[language=bash]
mkdir $HOME/cloudmesh
\end{lstlisting}
%$

We also assume you have installed gitbash and use the a git bash
terminal window for the next steps.

To install Docker on Windows, please download the following files

\URL{https://download.docker.com/win/stable/Docker\%20for\%20Windows\%20Installer.exe}

\URL{https://download.docker.com/win/stable/DockerToolbox.exe}

Move the downloaded files to a directory \verb|$HOME/cloudmesh|

Place the downloaded exe files in the cloudmesh directory we created
earlier. First conduct the Docker installation and then continue with
the Docker Toolbox installation. You can double click the exe files  and
run the installation.

When you are doing the installation, check on the selection in the
installation saying to create shortcuts in your desktop. This way you
can start all the programs from a desktop shortcut.

To run docker, you need to start the \textit{Docker Quickstart
  Terminal} application and it will load all docker and provide a
terminal window in which you can execute docker commands. Once the
terminal is loaded, it will show something like following:

\begin{lstlisting}[language=bash]
$ <username>@<yourpc> ~
\end{lstlisting}
%$



\section{Testing the Install}
\index{Docker!Install!Test}

To test if it works execute the following commands in a terminal:

\begin{lstlisting}[language=bash]
docker version
\end{lstlisting}

You should see an output similar to

\begin{lstlisting}
docker version

Client:
  Version:      17.03.1-ce
  API version:  1.27
  Go version:   go1.7.5
  Git commit:   c6d412e
  Built:        Tue Mar 28 00:40:02 2017
  OS/Arch:      darwin/amd64

Server:
  Version:      17.03.1-ce
  API version:  1.27 (minimum version 1.12)
  Go version:   go1.7.5
  Git commit:   c6d412e
  Built:        Fri Mar 24 00:00:50 2017
  OS/Arch:      linux/amd64
  Experimental: true
\end{lstlisting}

To see if you can run a container use

\begin{lstlisting}[language=bash]
docker run hello-world
\end{lstlisting}

Once executed you should see an output similar to

\begin{lstlisting}[language=bash]
Unable to find image 'hello-world:latest' locally
latest: Pulling from library/hello-world
78445dd45222: Pull complete 
Digest: sha256:c5515758d4c5e1e838e9cd307f6c6a .....
Status: Downloaded newer image for hello-world:latest

Hello from Docker!
This message shows that your installation appears to 
be working correctly.

To generate this message, Docker took the following steps:
1. The Docker client contacted the Docker daemon.
2. The Docker daemon pulled the "hello-world" image 
   from the Docker Hub.
3. The Docker daemon created a new container from that 
   image which runs the executable that produces the 
   output you are currently reading.
4. The Docker daemon streamed that output to the Docker 
   client, which sent it to your terminal.

To try something more ambitious, you can run an Ubuntu container 
with:

$ docker run -it ubuntu bash

Share images, automate workflows, and more with a free Docker ID:
https://cloud.docker.com/

For more examples and ideas, visit:
https://docs.docker.com/engine/userguide/
\end{lstlisting}
%$
